\section{Threshold variations and the scalar Hessian}
\label{sec:web-threshold-variation}

This supplement retains the distinct information obtained by selecting
near-minimizers of the relative defect, rather than of $P$ itself.
All asymptotic statements concern the balanced family
\eqref{webvar:family}, and all selected representatives have $s=-4$.
Set
\begin{equation}\label{webvar:threshold-functional}
 \gamma_j=\frac{\alpha_j}{4(4-\alpha_j)}=O(h_j),\qquad
 F_j=R-\frac{S}{4-\alpha_j}=P-\gamma_jS.
\end{equation}
The coefficient $\gamma_j$ is unrelated to the blow-up number $k_j$.
For large $j$, $\alpha_j<1$ and $\inf\J=\alpha_j$, so
$F_j=R(\J-\alpha_j)/(4-\alpha_j)\ge0$ on every metric.
On scalar-$-4$ representatives, $\J=\E$.

\subsection{Fourth-order Taylor selection}

\begin{proposition}[Threshold selection]\label{webvar:threshold-selection}
For every $\varepsilon_j\downarrow0$, one can choose $s_{g_j}=-4$
with $\E(g_j)-\alpha_j\le\varepsilon_j^2h_j^2$ such that
\begin{equation}\label{webvar:threshold-budgets}
 R=8\pi^2A_j+O(Q_j),\quad V=2\pi^2A_j+O(Q_j),\quad P=O(Q_j),
\end{equation}
and, uniformly for $H\in\mathcal T(g_j,A_j)$,
\begin{equation}\label{webvar:threshold-variations}
 |DF_j[H]|\le C\varepsilon_jQ_j,\qquad
 D^2F_j[H,H]\ge-C\varepsilon_jQ_j.
\end{equation}
This sequence has $P(g_j)-p_j=O(A_jh_j^2)$; it need not have the
arbitrarily small $P$-error in Proposition~\ref{webvar:selection}.
\end{proposition}
\begin{proof}
The definition of $\alpha_j$ and the exact conformal comparison allow
the prescribed defect accuracy on scalar-$-4$ representatives.
The two-sided curvature budget gives
\eqref{webvar:threshold-budgets}.  In particular
\[
 0\le F_j(g_j)
 =\frac{R(g_j)}{4-\alpha_j}(\E(g_j)-\alpha_j)
 \le C A_j\varepsilon_j^2h_j^2.
\]
The connection-difference calculation in
Lemma~\ref{webvar:variation-remainder} applies through fourth order:
differentiating in the affine metric parameter adds inverse metrics
and $H$, but no spatial derivatives beyond $\nabla^2H$.
For either $P$ or $S$, the second and fourth derivatives are bounded
by
\[
 C\int(|\Ric_{g_j}|^2+|\nabla^2H|^2+|\nabla H|^4)\le C A_j.
\]
The Ricci difference formula suffices here; no derivative of curvature
is required.  Since $\gamma_j$ is bounded, the same estimate holds
for $F_j(g_j+tH)$ on $|t|\le1/4$.

Let $f(t)=F_j(g_j+tH)$ and $d_j=C A_j\varepsilon_j^2h_j^2$.
Nonnegativity of $f(\pm t)$ and $f(0)\le d_j$ give
$|f'(0)|\le d_j/t+C A_jt$.  Take $t=\varepsilon_jh_j$.
The symmetric fourth-order expansion gives
$f''(0)\ge-2d_j/t^2-C A_jt^2$; take
$t=\sqrt{\varepsilon_jh_j}$.  These prove
\eqref{webvar:threshold-variations}.  Finally, on these representatives,
$\E=4P/(8\pi^2A_j)+O(h_j^2)$.
Comparison with \eqref{webvar:reduction} gives
$P(g_j)-p_j=O(A_jh_j^2)$.
\end{proof}

Write $Z=\Ric+g=\Ric^\circ$.  Since $DS[H]=8\int\langle Z,H\rangle$,
\begin{equation}\label{webvar:threshold-first-correction}
 |DP[H]|\le C\varepsilon_jQ_j+
       8\gamma_j\sqrt P\,\|H\|_2
       \le C(\varepsilon_j+\sqrt{h_j})Q_j=o(Q_j).
\end{equation}
Here $\|H\|_2\le\sqrt V$.  At second order the exact relation is
\begin{equation}\label{webvar:threshold-second-correction}
 D^2P[H,H]=D^2F_j[H,H]+\gamma_jD^2S[H,H].
\end{equation}
The general test budget only gives $|D^2S|=O(A_j)$, so this
correction can be $O(Q_j)$.  It cannot be discarded at first order
in the surgery cost.

\subsection{An exact identity without a TT projection}

Use the positive rough Laplacian $L=\nabla^*\nabla$ and the conventions
\[
 \Rm(H,K)=R_{ikjl}H^{ij}K^{kl},\qquad
 C(H)_{ijk}=\nabla_kH_{ij}-\nabla_jH_{ik}.
\]
The norm of $C(H)$ includes all ordered indices.  For trace-free $H$ set
$b_i=\nabla^jH_{ij}$, $\phi=\langle Z,H\rangle$, and
$\psi=\nabla^ib_i-\phi=Ds[H]$.  Also write
$Z:H^2=Z_{ij}H^i{}_kH^{kj}$.

\begin{proposition}[Trace-free scalar Hessian]
\label{webvar:scalar-hessian}
On any smooth closed four-manifold with $s=-4$, every smooth
trace-free symmetric $H$ satisfies
\begin{equation}\label{webvar:scalar-hessian-formula}
 \begin{split}
 D^2S[H,H]={}&2\|\psi\|_2^2+2\|C(H)\|_2^2-4\|b\|_2^2
              +\frac83\|H\|_2^2\\
 &-4\int W(H,H)-8\int Z:H^2 .
 \end{split}
\end{equation}
The derivative is along $g+tH$ and does not assume criticality.
For a TT tensor, $b=0$ and $\psi=-\phi$.
\end{proposition}
\begin{proof}
Put $I(g)=\int s_g$ and $\mathcal H(g)=I(g)+2V(g)$.
For every metric,
$S=\int(s+4)^2-8\mathcal H$.  At $s=-4$ this gives
\[
 D^2S[H,H]=2\|Ds[H]\|_2^2-8D^2\mathcal H[H,H].
\]
The covariant tensor gradient of $\mathcal H$ is
$-\Ric+(s/2+1)g=-Z$ at the base metric.  Differentiate this gradient,
the inverse metrics in its pairing with $H$, and the volume.
The last contribution vanishes because $\tr H=0$, and the result is
\[
 D^2\mathcal H[H,H]
 =-\int\langle D\Ric[H],H\rangle-\|H\|_2^2+2\int Z:H^2.
\]
For trace-free $H$, the Ricci linearization is
\[
 D\Ric[H]=\tfrac12\{LH+\Ric\circ H+H\circ\Ric-2\Rm(H)\}
                 +\operatorname{sym}\nabla b .
\]
Its pairing gives
\begin{equation}\label{webvar:hilbert-hessian}
 D^2\mathcal H[H,H]
 =-\tfrac12\|\nabla H\|_2^2+\int\Rm(H,H)+\|b\|_2^2+\int Z:H^2.
\end{equation}
One integration by parts and commutation of covariant derivatives give
\begin{equation}\label{webvar:cotton-integration}
 \tfrac12\|C(H)\|_2^2
 =\|\nabla H\|_2^2-\|b\|_2^2-\int\Rm(H,H)+\int\Ric:H^2.
\end{equation}
For example, expand the square defining $C(H)$; in its cross term,
move $\nabla_k$ to the other factor, commute it past $\nabla_j$,
and use $\nabla^kH_{ik}=b_i$.  The divergence term gives $\|b\|^2$,
and the two curvature actions give the last two terms.
The four-dimensional curvature decomposition on trace-free tensors is
\begin{equation}\label{webvar:curvature-on-tracefree}
 \Rm(H,H)=W(H,H)+\tfrac13|H|^2-Z:H^2.
\end{equation}
Substituting \eqref{webvar:cotton-integration} and
\eqref{webvar:curvature-on-tracefree} into
\eqref{webvar:hilbert-hessian} yields
\[
 D^2\mathcal H
 =-\tfrac14\|C(H)\|_2^2+\tfrac12\|b\|_2^2
      +\tfrac12\int W(H,H)-\tfrac13\|H\|_2^2+\int Z:H^2.
\]
Multiplication by $-8$ proves the formula.
\end{proof}

Several backgrounds check the coefficients independently.  For sectional
curvature $-1/3$, $Z=W=0$, and the TT formula becomes
$D^2S=2\|C(H)\|_2^2+(8/3)\|H\|_2^2$.
On a product of curvature-$-1$ surfaces, a Teichm\"uller tensor on
one factor is TT and Codazzi, and $W(H,H)=2|H|^2/3$, giving zero.
For the parallel TT direction $H=g_1-g_2$ on the same product,
$|H|^2=4$ and $W(H,H)=-16/3$; the formula gives $32V$.
Directly,
$s(g+tH)=-4/(1-t^2)$, $d\mu_{g+tH}=(1-t^2)d\mu_g$, and
$S(g+tH)=16V/(1-t^2)$, whose second derivative is also $32V$.

For a non-Einstein product with factor curvatures
$\kappa_1+\kappa_2=-2$, put $a=\kappa_1+1$ and take a
Teichm\"uller direction $H$ on the first factor.  It is Codazzi,
$Z:H^2=a|H|^2$, and $W(H,H)=2|H|^2/3$.
Thus $D^2S[H,H]=-8a\|H\|_2^2$.
Along the nonlinear Teichm\"uller path, $S$ is constant and the
fixed-area constraint gives
$\int\tr(g_1'')=\int|H|^2$.
Since $DS[g_1'']=8a\int\tr(g_1'')$, differentiation of
$S(g_t)=\mathrm{constant}$ confirms the same affine Hessian.
This check retains the inverse-metric derivative responsible for
the coefficient $-8$.

\subsection{The actual Weyl budget and threshold stability}

Let $w_3^\pm$ be the largest eigenvalues of the two Weyl operators and set
$\xi_W=(w_3^++w_3^--1/3)_+$.  On the late balanced family
$U_j=3T_-/A_j$.  For the actual class with $p\ge A_j$, put
\begin{equation}\label{webvar:weyl-budget}
 \begin{split}
 \mathcal B_j&=D_2/18+3W_-\\
 &=P+2\pi^2A_j\{U_j-1-h_j-2(p/A_j-1)\}
 \le P+O(A_jh_j^2)=O(Q_j).
 \end{split}
\end{equation}
The equality follows by substituting the two chiral identities
and $D_2=24V+18W_+-72\pi^2p$.
The model $W_0$ in \eqref{webvar:single-decomposition} has largest
eigenvalue $\rho^2/12\le1/3$.
For a trace-free three-dimensional symmetric operator, the largest
eigenvalue is at most $\sqrt{2/3}$ times its Hilbert--Schmidt norm.
Consequently
\[
 \xi_W\le\sqrt{2/3}(|W^+-W_0|+|W^-|),\qquad
 \int\xi_W^2\le\frac89(D_2/18+3W_-)=\frac89\mathcal B_j.
\]
Here $(a+b)^2\le(4/3)(a^2+3b^2)$ and the exact field decomposition
control the last step.  Under the standard algebraic identification
$S_0^2T^*\cong\Lambda^+\otimes\Lambda^-$,
\[
 W(H,H)\le(w_3^++w_3^-)|H|^2\le(1/3+\xi_W)|H|^2.
\]
For $\|H\|_\infty\le1$, H\"older's inequality therefore gives
\[
 \int(4\xi_W+8|Z|)|H|^2
 \le(4\|\xi_W\|_2+8\|Z\|_2)\|H\|_4^2
 \le C\sqrt{A_jQ_j}.
\]
Combine this with \eqref{webvar:scalar-hessian-formula} and
\eqref{webvar:threshold-second-correction} to obtain
\begin{equation}\label{webvar:tracefree-threshold-stability}
 \begin{split}
 D^2P[H,H]\ge{}&
 \gamma_j\{2\|Ds[H]\|_2^2+2\|C(H)\|_2^2
                  +\tfrac43\|H\|_2^2-4\|\operatorname{div}H\|_2^2\}\\
 &-C(\varepsilon_j+\sqrt{h_j})Q_j .
 \end{split}
\end{equation}
This holds uniformly for all trace-free $H\in\mathcal T(g_j,A_j)$.
For TT tests the divergence penalty vanishes and $Ds[H]=-\langle Z,H\rangle$.
More generally $\|\operatorname{div}H\|_2^2=o(A_j)$ makes its penalty
$o(Q_j)$, without solving a TT projection problem.

Along a subsequence with $\alpha_j/h_j\to c\in(0,4]$ one has
$\gamma_j/h_j\to c/16$.  Thus for the TT tests,
\begin{equation}\label{webvar:threshold-spectral-constraint}
 \frac{D^2P[H,H]}{Q_j}\ge
 \left(\frac c8+o(1)\right)
   \frac{\|C(H)\|_2^2+\|\langle Z,H\rangle\|_2^2}{A_j}
 +\left(\frac c{12}+o(1)\right)\frac{\|H\|_2^2}{A_j}-o(1).
\end{equation}
The displayed norm ratios are bounded by the amplitude, critical
budget, and $P=O(Q_j)$.  This necessary positive second-order cost
does not imply $c\ge4$: a smaller $c$ demands a smaller cost.
No admissible TT direction contradicting this inequality has been
constructed here.  Although $Z$ is TT at constant scalar curvature,
its amplitude and critical derivative budget are not controlled,
so the substitution $H=Z$ is not licensed by this selection.

\section{Adaptive variations and positive curvature compensation}
\label{sec:web-curvature-compensation}

The following tests are bounded without truncating a concentrated
curvature source before inversion.  They are admissible nonconformal
variations, but admissibility alone does not determine their first
or second variation.

\subsection{An adaptive rough-Laplacian resolvent}

\begin{lemma}[A bounded adaptive test]\label{webvar:adaptive-test}
Let $g$ be a smooth metric on a closed four-manifold with
$\int|\Rm|^2=O(A)$.  Let $B$ be a smooth symmetric trace-free
two-tensor, and let $\Upsilon$ be a smooth tensor satisfying
$|B|\le c_*|\Upsilon|$ and $\int|\Upsilon|^2=O(Q)$.
Set
\begin{equation}\label{webvar:adaptive-equation}
 v=\sqrt{1+|\Upsilon|^2},\qquad (L+v)H=B,\qquad L=\nabla^*\nabla.
\end{equation}
There is a unique smooth solution in the symmetric trace-free bundle,
and it satisfies
\begin{equation}\label{webvar:adaptive-graph}
 \|H\|_\infty\le c_*,\qquad
 \int(|H|^2+|\nabla H|^2+|LH|^2)=O(Q),
\end{equation}
\begin{equation}\label{webvar:adaptive-critical-general}
 \int(|\nabla^2H|^2+|\nabla H|^4)=O(Q+A).
\end{equation}
If $Q=O(A)$ the last bound is $O(A)$.  Constants depend on the
tensor type and stated budget constants, not on geometric compactness
or derivatives of $\Upsilon$.  On the actual balanced family,
\begin{equation}\label{webvar:adaptive-critical-sharp}
 \int(|\nabla^2H|^2+|\nabla H|^4)=O(Q).
\end{equation}
\end{lemma}
\begin{proof}
The positive potential $v\ge1$ makes $L+v$ invertible; the induced
connection preserves symmetry and trace.  Kato's inequality and
the maximum principle, or their regularized form at zeros, give
$(\Delta+v)|H|\le|B|\le c_*v$, so $\|H\|_\infty\le c_*$.
Pairing \eqref{webvar:adaptive-equation} with $H$ gives
\[
 \int(|\nabla H|^2+v|H|^2)=\int\langle B,H\rangle .
\]
Hence $\|H\|_2\le\|B\|_2=O(\sqrt Q)$ and
$\|\nabla H\|_2^2=O(Q)$.
Moreover
$\int v^2|H|^2=\int|H|^2+\int|\Upsilon|^2|H|^2=O(Q)$,
so $LH=B-vH$ proves the remaining graph-norm bound.

The required Hessian estimate can be obtained without differentiating
$v$.  Integration of a divergence gives
\begin{equation}\label{webvar:bounded-section-interpolation}
 \int|\nabla H|^4\le C\|H\|_\infty^2\int|\nabla^2H|^2.
\end{equation}
Indeed, integrate one derivative in
$\nabla H*\nabla H*\nabla H*\nabla H$ onto the other factors.
The result is bounded by
$C\|H\|_\infty\int|\nabla^2H||\nabla H|^2$,
and Cauchy--Schwarz proves the assertion.
Commuting derivatives in the integrated Hessian square and moving
curvature derivatives off the curvature yields terms of the forms
$\Rm*H*\nabla^2H$ and $\Rm*\nabla H*\nabla H$, besides $\|LH\|_2^2$.
H\"older, \eqref{webvar:bounded-section-interpolation}, and Young give
\[
 \int|\nabla^2H|^2
 \le C\left(\int|LH|^2+\|H\|_\infty^2\int|\Rm|^2\right).
\]
Together with the graph bound this proves
\eqref{webvar:adaptive-critical-general}.

For the actual balanced family, \eqref{webvar:weyl-budget} gives
$\|\xi_W\|_2^2=O(Q)$.  A trace-free three-dimensional Weyl block has
nonnegative largest eigenvalue, and all its eigenvalues have absolute
value at most twice that largest eigenvalue.  Thus, at $s=-4$,
\begin{equation}\label{webvar:actual-curvature-envelope}
 |\Rm|\le C(1+e),\qquad e=\xi_W+|Z|,\qquad \|e\|_2^2=O(Q).
\end{equation}
In the Hessian commutation estimate, the resulting upper bound is
\[
 \|LH\|_2^2+C\int(1+e)|\nabla H|^2
 +C\int(1+e)|H||\nabla^2H|+C\int(1+e)^2|H|^2.
\]
Set $X=\|\nabla^2H\|_2^2$.
The bounded-coefficient terms use \eqref{webvar:adaptive-graph}.
The remaining terms are at most
$C Q+C\sqrt Q\,\sqrt X$, by $\|H\|_\infty\le c_*$,
$\|e\|_2=O(\sqrt Q)$, and
$\|\nabla H\|_4^2\le C\|H\|_\infty\sqrt X$.
Young's inequality absorbs $\sqrt X$, yielding $X=O(Q)$.
Equation~\eqref{webvar:bounded-section-interpolation} proves
\eqref{webvar:adaptive-critical-sharp}.  No derivative of $e$,
$\Upsilon$, or the Weyl eigenvalues has been used.
\end{proof}

After a fixed rescaling, these tensors belong to
$\mathcal T(g_j,A_j)$.  They also satisfy
$\|\operatorname{div}H\|_2^2=O(Q_j)=o(A_j)$.
The exact scalar-Hessian formula
\eqref{webvar:scalar-hessian-formula} now gives
\begin{equation}\label{webvar:adaptive-scalar-correction}
 |D^2S[H,H]|=O(Q_j),\qquad \gamma_jD^2S[H,H]=o(Q_j).
\end{equation}
To see this directly, $Ds[H]=\operatorname{div}\operatorname{div}H-
\langle Z,H\rangle$ has squared $L^2$ norm $O(Q_j)$; the Cotton,
divergence, and tensor norms have the same bound.  The remaining
curvature terms use \eqref{webvar:actual-curvature-envelope} and
$\|H\|_4^2\le\|H\|_\infty\|H\|_2$.
Thus the second-order threshold correction is lower order for these
localized tests.  This differs from a bulk test with mass and
derivative budget of order $A_j$, for which the correction can be
of order $Q_j$.

Two concrete sources will be useful.  The direct Ricci test is
\begin{equation}\label{webvar:direct-resolvent}
 (L+\sqrt{1+|Z|^2})H_Z=Z,\qquad \|H_Z\|_\infty\le1,
\end{equation}
and it has the positive pairing
\begin{equation}\label{webvar:direct-pairing}
 \int\langle Z,H_Z\rangle
  =\int(|\nabla H_Z|^2+\sqrt{1+|Z|^2}|H_Z|^2)\ge0.
\end{equation}
For the same ordinary solution as in
\eqref{webvar:ordinary-fields}, put
$\omega=iF_{\mathcal A}^+$ and $\eta=iF_{\mathcal A}^-$ and define
\begin{equation}\label{webvar:mixed-resolvent}
 B_{ij}=\omega_i{}^k\eta_{jk}+\eta_i{}^k\omega_{jk},\qquad
 (L+\sqrt{1+|\eta|^2})H_{\mathcal A}=B.
\end{equation}
Opposite chirality makes $B$ symmetric and trace-free.
The spinor maximum principle gives $|\omega|\le\sqrt2$, so
$|B|\le c_*|\eta|$.  Also
\[
 \int|\eta|^2=\int|\omega|^2-4\pi^2a^2
 \le2V-4\pi^2(A-k)=O(Q).
\]
Lemma~\ref{webvar:adaptive-test} therefore applies without truncating
$B$ or differentiating $\eta$.

The potential retains the source scale.  In a Euclidean local model,
write $B_\ell(x)=\ell^{-2}\overline B(x/\ell)$ and
$\Upsilon_\ell(x)=\ell^{-2}\overline\Upsilon(x/\ell)$.
With $x=\ell y$, the equation becomes
\[
 \left(-\sum_a\partial_{y_a}^2+
       \sqrt{\ell^4+|\overline\Upsilon|^2}\right)\overline H_\ell
       =\overline B.
\]
There is no residual factor $\ell^2$ on the source, so an order-one
test amplitude and a scale-invariant Hessian budget are possible.
Truncating the source to bounded amplitude before inversion at unit
scale would leave such a factor and diminish the core test.
This calculation is a scaling model, not a newly constructed monopole
core.  In \eqref{webvar:mixed-resolvent}, vanishing of $\omega$ can
still make the mixed source fail to detect a core.

\subsection{A positive nonlinear identity}

\begin{proposition}[Resolvent compensation]\label{webvar:resolvent-positive}
Let $Z$ be any smooth section of a Euclidean vector bundle with a
metric connection over a smooth closed Riemannian manifold.
Set $v=\sqrt{1+|Z|^2}$ and $(L+v)H=Z$.
Then
\begin{equation}\label{webvar:positive-resolvent}
 \mathfrak D(Z):=\int|Z|^2-\int v\langle Z,H\rangle
 \ge\int|\nabla H|^2+\int v^{-1}|LH|^2\ge0.
\end{equation}
The statement is independent of dimension and curvature sign.
\end{proposition}
\begin{proof}
Let $z=Z/v$ and define the fiberwise map
$\mathcal F(u)=u/\sqrt{1-|u|^2}$ on the open unit ball.
Compactness gives
$m=\max_M|Z|/\sqrt{1+|Z|^2}<1$.
From $(L+v)H=vz$ and the maximum principle, $|H|\le m<1$.
Thus $\mathcal F(H)$ is a smooth global section; no uniform
separation of $m$ from one across a sequence is needed.
One has $\mathcal F(z)=Z$ and $LH=v(z-H)$.
The fiber derivative is
\[
 D\mathcal F(u)[w]
 =\frac{w}{\sqrt{1-|u|^2}}
   +\frac{u\langle u,w\rangle}{(1-|u|^2)^{3/2}}.
\]
Hence $\langle D\mathcal F(u)[w],w\rangle\ge|w|^2$.
Integration along a segment in the unit ball gives
$\langle\mathcal F(u)-\mathcal F(w),u-w\rangle\ge|u-w|^2$.
The exact decomposition is
\begin{equation}\label{webvar:positive-resolvent-decomposition}
 \begin{split}
 \mathfrak D(Z)
 &=\int v\langle\mathcal F(z)-\mathcal F(H),z-H\rangle
          +\int\langle\mathcal F(H),LH\rangle\\
 &=\int v\langle\mathcal F(z)-\mathcal F(H),z-H\rangle
          +\int\langle D\mathcal F(H)[\nabla H],\nabla H\rangle.
 \end{split}
\end{equation}
Metric compatibility gives
$\nabla\mathcal F(H)=D\mathcal F(H)[\nabla H]$.
Applying both monotonicity inequalities and using $LH=v(z-H)$
proves \eqref{webvar:positive-resolvent}.
\end{proof}

The same proof with
$v_\lambda=\sqrt{\lambda^2+|Z|^2}$, $\lambda>0$,
$(L+v_\lambda)H_\lambda=Z$, and
$\mathcal F_\lambda(u)=\lambda u/\sqrt{1-|u|^2}$ gives
\begin{equation}\label{webvar:positive-resolvent-lambda}
 \int|Z|^2-\int v_\lambda\langle Z,H_\lambda\rangle
 \ge\lambda\int|\nabla H_\lambda|^2
       +\lambda\int v_\lambda^{-1}|LH_\lambda|^2.
\end{equation}
The positivity concerns the difference of the two integrals in
\eqref{webvar:positive-resolvent}; it does not assert positivity of
the weighted pairing alone.

\subsection{First variations and the remaining curvature term}

At $s=-4$, the classical first-variation formulas are
\begin{equation}\label{webvar:curvature-gradients}
 \operatorname{grad}P=LZ-2\Rm(Z)+\tfrac12|Z|^2g,\qquad
 \operatorname{grad}S=8Z.
\end{equation}
They follow from the Ricci-square and scalar-square formulas in
\cite[Section~2.1, equations (2.5) and (2.9)]{GurskyViaclovsky},
with the positive rough-Laplacian convention and $P=R-S/4$.
For fixed $\gamma$ and any trace-free $H$,
\begin{equation}\label{webvar:resolvent-first-general}
 D(P-\gamma S)[H]
 =\int\langle Z,LH-2\Rm(H)-8\gamma H\rangle.
\end{equation}
Self-adjointness of the curvature action and integration by parts
justify the expression even when $\operatorname{div}H\ne0$.
For $H_Z$ from \eqref{webvar:direct-resolvent}, it becomes
\begin{equation}\label{webvar:resolvent-first-direct}
 D(P-\gamma S)[H_Z]
 =P-\int(v_Z+8\gamma)\langle Z,H_Z\rangle
                  -2\int\Rm(Z,H_Z),\quad v_Z=\sqrt{1+|Z|^2}.
\end{equation}
For \eqref{webvar:mixed-resolvent} it becomes
\begin{equation}\label{webvar:resolvent-first-mixed}
 D(P-\gamma S)[H_{\mathcal A}]
 =\int\langle Z,B\rangle-\int(v_{\mathcal A}+8\gamma)
       \langle Z,H_{\mathcal A}\rangle-2\int\Rm(Z,H_{\mathcal A}).
\end{equation}
Every tensor in the latter equation belongs to the same metric and
the same actual monopole solution.

For the $P$-selection in Proposition~\ref{webvar:selection},
\eqref{webvar:resolvent-first-direct} and
\eqref{webvar:positive-resolvent} give
\begin{equation}\label{webvar:forced-curvature-compensation}
 2\int\Rm(Z,H_Z)
 =\mathfrak D(Z)+o(Q)
 \ge\int|\nabla H_Z|^2+\int v_Z^{-1}|LH_Z|^2-o(Q).
\end{equation}
For the threshold selection the same conclusion holds, because
$\gamma=O(h)$ and $\int\langle Z,H_Z\rangle=O(Q)$.
This uses the appropriate weak stationarity separately on each
selection; it does not identify the two minimizing sequences.
Equivalently,
\[
 P=\int v_Z\langle Z,H_Z\rangle+2\int\Rm(Z,H_Z)+o(Q).
\]
The graph estimate implies
$|\int v_Z\langle Z,H_Z\rangle|
\le\|Z\|_2\|v_ZH_Z\|_2=O(Q)$ and $\mathfrak D(Z)=O(Q)$.
The actual curvature envelope also gives, independently of stationarity,
\begin{equation}\label{webvar:curvature-coupling-size}
 \int|\Rm(Z,H)|
 \le C\|Z\|_2\|H\|_2+
       C\|H\|_\infty\|\xi_W+|Z|\|_2\|Z\|_2=O(Q).
\end{equation}
This improves the $O(\sqrt{AQ})$ bound from the full curvature
budget alone, but is not $o(Q)$.  Equation
\eqref{webvar:forced-curvature-compensation} forces the curvature
coupling to pay a nonnegative quantity; it contains no term
$8\pi^2Q$ and does not prove the sharp topological cost.

All these variations hold with $H$ fixed after solving its equation
at the base metric.  Differentiating the displayed resolvent identities
as functions of the base metric introduces $D_gH$ and $D_gv$.
Such a derivative is not automatically the affine second variation.

\subsection{The fourth-order symbol and the volume route}

The correction $-\gamma S$ does not produce arbitrary high-frequency
negative directions.  At a point choose a frequency
$\zeta=|\zeta|e^0$ and, modulo infinitesimal diffeomorphisms, use
$H_{0i}=0$.  Write its spatial block as
$H_{ab}=K_{ab}+(t/3)\delta_{ab}$, $\tr K=0$.
The principal Ricci symbols have entries
$\sigma(D\Ric)_{00}=|\zeta|^2t/2$ and
$\sigma(D\Ric)_{ab}=|\zeta|^2H_{ab}/2$, up to a common sign.
Thus
\[
 |\sigma(D\Ric^\circ)H|^2
 =|\zeta|^4(|K|^2/4+t^2/12),\qquad
 |\sigma(Ds)H|^2=|\zeta|^4t^2.
\]
For $\gamma=\alpha/[4(4-\alpha)]$ the fourth-order symbol of the
Hessian is consequently
\begin{equation}\label{webvar:threshold-principal-symbol}
 |\zeta|^4\left\{\frac12|K|^2+
                    \frac{2(1-\alpha)}{3(4-\alpha)}t^2\right\}.
\end{equation}
It is positive for $\alpha<1$ and uniformly positive as $\alpha\to0$.
On each fixed smooth background, frequencies sufficiently above its
curvature and variation scales therefore have a dominating positive
fourth-order term.  This does not exclude a test at the curvature
concentration scale, or a macroscopic or low-frequency variation.

There is a distinct sufficient scalar condition.  Keep separate
$\Delta_S=S-32\pi^2A=16(V-2\pi^2A)$ and
$\Sigma=S-32\pi^2p$, where $p\ge A$ is the actual projection square.
The classical Ricci estimate gives
\begin{equation}\label{webvar:scalar-slack-route}
 P\ge8\pi^2Q+8\pi^2(p-A)-\Sigma/4.
\end{equation}
Indeed, subtract $S/4$ from $R\ge8\pi^2(2p-T_+)$ and use
$T_+=A-Q$.  Thus $\Sigma=o(Q)$ suffices for coefficient four.
The stronger condition $\Delta_S=o(Q)$, equivalently
$V=2\pi^2A+o(Q)$, also suffices, because
$0\le\Sigma\le\Delta_S$.

The degree-one map $M\to X$ allows the Besson--Courtois--Gallot
entropy inequality \cite{BessonCourtoisGallot1995},
\begin{equation}\label{webvar:entropy-volume}
 h_{\rm vol}(g)^4V(g)\ge h_{\rm vol}(g_X)^4V(g_X).
\end{equation}
It is a lower bound on a product.  An entropy upper bound still
gives a volume lower bound, whereas the scalar route requires an
upper bound close to the known minimum.  Even under the additional
pointwise hypothesis $\Ric_g\ge-g$, Bishop comparison gives
$h_{\rm vol}(g)^2\le3$.  In the normalization $\Ric_{g_X}=-g_X$,
the complex-hyperbolic entropy has square $8/3$.  Hence this
combination gives only $V(g)\ge(64/81)V(g_X)$, weaker than
the monopole bound $V(g)\ge V(g_X)$.  The sharp Ricci-volume
corollary in \cite{BessonCourtoisGallot1995} concerns real-hyperbolic
targets and cannot be transferred in that form to this target.

The entropy estimate of \cite[Theorem~1.4]{ChenWei} requires a
small integral negative-Ricci error with exponent $p>n/2$ and
diameter control.  The present four-dimensional critical $p=2$
budget supplies neither hypothesis.  Its real-hyperbolic stability
theorem additionally assumes a nearly minimal volume upper bound
\cite[Theorem~1.5]{ChenWei}; that conclusion cannot manufacture
the same bound as a premise.  These observations limit this direct
use of entropy comparison and do not exclude a new coupled
entropy-curvature theorem.

\section{The Codazzi decomposition and its scope}
\label{sec:web-codazzi-scope}

The positive rough-resolvent identity does not transfer automatically
to another elliptic system.  We give the exact bilinear Codazzi
decomposition, identify its symbol, and exhibit the pointwise failure
of the radial positivity step that such a transfer would need.
All formulas use the conventions of
Section~\ref{sec:web-threshold-variation}, with $s=-4$ and $Z=\Ric+g$.
The contracted Bianchi identity gives $\operatorname{div}Z=0$.

\subsection{Bilinear decomposition and principal symbol}

\begin{proposition}[Codazzi first variation]
\label{webvar:codazzi-first-variation}
Every smooth symmetric trace-free $H$, without a divergence condition,
satisfies
\begin{equation}\label{webvar:codazzi-P}
 DP[H]=\frac12\int\langle C(Z),C(H)\rangle
             +\int\left\langle\left(\frac23\Id-W\right)Z,H\right\rangle.
\end{equation}
For fixed $\gamma$,
\begin{equation}\label{webvar:codazzi-threshold}
 D(P-\gamma S)[H]=\frac12\int\langle C(Z),C(H)\rangle
       +\int\left\langle
              \left(\left(\frac23-8\gamma\right)\Id-W\right)Z,H\right\rangle.
\end{equation}
\end{proposition}
\begin{proof}
Polarize \eqref{webvar:cotton-integration}.  Since
$\operatorname{div}Z=0$, the divergence pairing vanishes, and
\[
 \frac12\int\langle C(Z),C(H)\rangle
 =\int\langle\nabla Z,\nabla H\rangle
             -\int\Rm(Z,H)+\int\Ric:ZH.
\]
Here $\Ric:ZH=-\langle Z,H\rangle+\tr(Z^2H)$.
The curvature decomposition is
\[
 \Rm(Z,H)=W(Z,H)+\frac13\langle Z,H\rangle-\tr(Z^2H).
\]
Solve the first equation for the derivative pairing and substitute
into the trace-free test of \eqref{webvar:curvature-gradients},
$DP[H]=\int\langle\nabla Z,\nabla H\rangle-2\int\Rm(Z,H)$.
The cubic traces cancel, giving \eqref{webvar:codazzi-P}.
Subtracting $\gamma DS[H]=8\gamma\int\langle Z,H\rangle$
gives \eqref{webvar:codazzi-threshold}.
\end{proof}

For a coefficient check, take the product of surfaces with curvatures
$-1+a$ and $-1-a$.  Then $Z=a(g_1-g_2)$ is parallel and
$W(Z,Z)=-(4/3)|Z|^2$, so \eqref{webvar:codazzi-P} gives $DP[Z]=2P$.
Along the affine path,
\[
 P(g+tZ)=\frac{4a^2(1+t)^2}{1-a^2t^2}\,V(g),
\]
which has the same derivative at zero.

The distributional operator in \eqref{webvar:codazzi-threshold}
on the trace-free tensor bundle is
\[
 K_\gamma=\tfrac12C^*C+
          \left(\left(\tfrac23-8\gamma\right)\Id-W\right).
\]
For a unit covector $\zeta$, its principal quadratic form is
\[
 \tfrac12|\sigma_C(\zeta)H|^2=|H|^2-|H\zeta|^2.
\]
To estimate it, take $\zeta=e_1$ and write $H$ in blocks with
first diagonal entry $a$, remaining first column $b$, and remaining
$3$ by $3$ block $B$.  Trace-freeness gives $\tr B=-a$, whence
$|B|^2\ge a^2/3$.  Therefore
$|H\zeta|^2=a^2+|b|^2\le(3/4)|H|^2$, and
\begin{equation}\label{webvar:codazzi-symbol}
 \frac14|H|^2\le\tfrac12|\sigma_C(\zeta)H|^2\le|H|^2.
\end{equation}
Thus the system is strongly elliptic.  Its zero-order form satisfies
\begin{equation}\label{webvar:codazzi-potential}
 \left\langle
    \left(\left(\tfrac23-8\gamma\right)\Id-W\right)H,H\right\rangle
 \ge(\tfrac13-8\gamma-\xi_W)|H|^2.
\end{equation}
Strong ellipticity is not Kato domination or a scalar maximum principle.
It does not supply the $L^\infty$ estimate enjoyed by
$(L+\sqrt{1+|Z|^2})^{-1}$.

\subsection{Failure of radial positivity for the Codazzi operator}

The proof of Proposition~\ref{webvar:resolvent-positive} uses
$D\mathcal F\ge\Id$ for
$\mathcal F(U)=U/\sqrt{1-|U|^2}$, which makes
$\langle\nabla\mathcal F(H),\nabla H\rangle$ nonnegative.
The corresponding pointwise statement for $C$ is false.
In Euclidean four-space, with $\zeta=e_1$, set
\[
 U=\frac r{\sqrt2}\operatorname{diag}(1,-1,0,0),\qquad
 V=\operatorname{diag}(2,1,-3,0),\qquad 0<r<1.
\]
Prescribe the first jet $H=U$, $\partial_1H=V$, and
$\partial_iH=0$ for $i\ne1$.  It is realized by a smooth local
trace-free tensor remaining in the unit ball on a sufficiently small
neighborhood.  If $f=(1-r^2)^{-1/2}$, then
$D\mathcal F(U)V=fV+f^3\langle U,V\rangle U$.
For the symbol pairing
$\mathcal B(A,B)=\langle A,B\rangle-\langle A\zeta,B\zeta\rangle$,
direct calculation gives
\[
 \mathcal B(V,V)=10,\qquad
 \langle U,V\rangle=r/\sqrt2,\qquad
 \mathcal B(U,V)=-r/\sqrt2.
\]
Hence
\begin{equation}\label{webvar:codazzi-radial-counterexample}
 \frac12\langle C(\mathcal F(H)),C(H)\rangle
       =10f-\frac{r^2}{2}f^3<0
       \quad\text{when }r^2>\frac{20}{21}.
\end{equation}
At $r^2=99/100$ the value is $-395$.
This refutes the pointwise positivity used in a direct substitution
of the Codazzi operator into the previous monotonicity proof.
It does not refute an independently compensated global integral
inequality or every possible Codazzi-resolvent method.

At a smooth critical point of $P-\gamma S$, the unrestricted test
$H=Z$ is legitimate.  Equations~\eqref{webvar:codazzi-threshold} and
\eqref{webvar:codazzi-potential} then yield
\begin{equation}\label{webvar:stationary-codazzi}
 \frac12\|C(Z)\|_2^2+
                 \left(\frac13-8\gamma\right)P
             \le\int\xi_W|Z|^2.
\end{equation}
This is the diagonal Cotton information already present in the
manuscript.  On the balanced family, if $P\asymp Q$ and
$\gamma=O(h)$, \eqref{webvar:weyl-budget} and Cauchy--Schwarz give
\begin{equation}\label{webvar:stationary-fourth-moment}
 \int|Z|^4\ge\left(\frac18-o(1)\right)P.
\end{equation}
Indeed, square
$(1/3-o(1))P\le\|\xi_W\|_2(\int|Z|^4)^{1/2}$
and use $\|\xi_W\|_2^2\le(8/9)\{P+O(Ah^2)\}$.
This precludes distributing all defect with amplitude $O(\sqrt h)$
over volume of order $A$.  A general selected near-minimizer only
has weak stationarity for the specified bounded test class, however;
it need not admit the test $H=Z$.  Consequently
\eqref{webvar:stationary-codazzi}--\eqref{webvar:stationary-fourth-moment}
cannot be transferred to that sequence without further bounds.
For a legal adaptive test, the bilinear Codazzi pairing itself has
no established sign, as \eqref{webvar:codazzi-radial-counterexample}
explains.

\subsection{Yamabe correction and the variation of the class budget}

Let $H$ be trace-free, $g_t=g+tH$, and
$\widehat g_t=u_t^2g_t$ its scalar-$-4$ representative, with $u_0=1$.
Differentiating the scalar transformation law gives
\begin{equation}\label{webvar:yamabe-variation}
 (6\Delta+8)\dot u=-Ds[H],\qquad
 Ds[H]=\operatorname{div}\operatorname{div}H-\langle Z,H\rangle,
\end{equation}
where $\Delta=d^*d$ on functions.  Integration on the closed manifold
therefore gives
\begin{equation}\label{webvar:yamabe-volume-variation}
 \dot V=4\int\dot u=\tfrac12\int\langle Z,H\rangle,\qquad
 \dot S=16\dot V=8\int\langle Z,H\rangle.
\end{equation}
At constant scalar curvature the purely conformal first variation of
$P=4W_++S/12-8\pi^2T_+$ is zero.  Thus
\[
 \left.\frac d{dt}P(\widehat g_t)\right|_{t=0}=DP[H].
\]
Moving to the Yamabe slice does not turn small $DP$ into small $DS$.
For the direct adaptive test,
$\int\langle Z,H_Z\rangle
=\int(|\nabla H_Z|^2+v_Z|H_Z|^2)\ge0$,
but the present bound is only $O(Q)$.

For a fixed actual class $a$, let $\vartheta$ be its real harmonic
representative, normalized by $[\vartheta]=a$, and set
$p=(a_g^+)^2$.  Its stress tensor is
\[
 T(\vartheta)_{ij}
 =\vartheta_{ik}\vartheta_j{}^k
             -\tfrac14\vartheta_{ab}\vartheta^{ab}g_{ij}.
\]
The variation of $p$ is
\begin{equation}\label{webvar:projection-variation}
 p'[H]=-\frac12\int\langle H,T(\vartheta)\rangle
       =-\int H^{ij}\vartheta^+_{ik}(\vartheta^-)_j{}^k.
\end{equation}
To verify it, use
$p=(\|\vartheta\|_2^2+a^2)/2$.
The change of harmonic representative is exact, so its pairing with
$\vartheta$ integrates to zero.  Variation of the inverse metrics
and volume in the remaining norm gives
$-\int\langle H,T(\vartheta)\rangle$.
Pure self-dual and anti-self-dual stresses vanish in four dimensions,
leaving the mixed expression in \eqref{webvar:projection-variation}.
Conformal invariance of the middle-degree Hodge star makes the formula
unchanged by the Yamabe correction.
Consequently the scalar slack satisfies
\begin{equation}\label{webvar:slack-variation}
 \Sigma_a'[H]
 =8\int\langle Z,H\rangle+16\pi^2\int\langle H,T(\vartheta)\rangle,
 \qquad \Sigma_a=S-32\pi^2p.
\end{equation}
Different actual sign classes give different harmonic stresses.
No proved identity makes the relevant stress equal to
$-Z/(2\pi^2)$ or makes their sum supply the missing sign.
An inequality with an unsaturated remainder of order $Q$ also
cannot be differentiated as though it were an equality at a
$P$-near-minimizer.

\subsection{A formal stable minimum compatible with the scalar budgets}

The following is a model of energy data, not a metric construction.
In the balanced first-order normalization put
$t=P/(8\pi^2Q)$,
$x=(S-32\pi^2A)/(32\pi^2Q)$, and $b=W_-/(2\pi^2Q)$.
For any fixed $0<\eta<1/2$, let
\[
 t(x)=1-2\eta x+\eta x^2,\qquad
 b(x)=t(x)+\frac13-\frac x3,\qquad 0\le x\le1+\delta,
\]
where $\delta>0$ is sufficiently small.  Then
$t+x-1=(1-2\eta)x+\eta x^2\ge0$ and $b>0$.
These are the sharp Ricci constraint and the opposite-chirality
nonnegativity constraint at this order.  With the formal choice $p=A$,
the corresponding identities are
\begin{equation}\label{webvar:formal-stable-budget}
 D_2=36\pi^2Q(x+t-1)\ge0,\qquad D_2/18+3W_-=P.
\end{equation}
For the second, substitute
$W_-=2\pi^2Q(t+1/3-x/3)$.
At the recovery endpoint $x=0$, $t=1$.
At the interior point $x=1$,
$t'=0$, $t''=2\eta>0$, and $t=1-\eta<1$, while
$S-32\pi^2A=32\pi^2Q$.
Thus the scalar budgets and a stable one-variable minimum can
coexist formally with a lower cost and scalar slack of order $Q$.
This supplies no actual metric or solution of the field equations.
It identifies the extra work a sharp proof must do: use geometric
or simultaneous-field information beyond these scalar constraints
and the sign of a formal variation.  The global topological
compensation, and hence the unrestricted O11 coefficient, remains open.
